% =====================================================================
\chapitre{Présentation de l'application réalisée}
% =====================================================================

Après la conception et la réalisation, ce dernier chapitre montre le résultat. Il présente
\equipay telle qu'elle a été testée sur iPhone, à travers un parcours complet : un voyage créé, un membre qui le rejoint, deux dépenses ajoutées avec l'IA, les
comptes équilibrés, puis le programme de la journée organisé. Toutes les captures sont
issues d'une utilisation réelle de l'application, dans la devise de l'activité (le dirham).

% ---------------------------------------------------------------------
\subsection{Aperçu général}

L'écran d'accueil (figure~\ref{fig:ecran_accueil}, à gauche) liste les activités de
l'utilisateur ; le bouton orange permet d'en créer une ou d'en rejoindre une avec son code,
et le mode « Sélectionner » permet d'archiver ou de supprimer plusieurs activités d'un coup.
La barre d'onglets donne accès aux activités, aux demandes pour rejoindre et au profil ; un
quatrième onglet, « Free Card », est réservé à une évolution future. Le profil regroupe les
informations du compte, les activités archivées et les préférences (langue, mode sombre). La
bulle « Support » ouvre l'assistant conversationnel depuis n'importe quel écran principal.

\begin{figure}[H]
\centering
\begin{subfigure}[t]{0.31\textwidth}\centering\capture[0.93]{accueil.jpeg}\caption{Accueil : mes activités}\end{subfigure}\hfill
\begin{subfigure}[t]{0.31\textwidth}\centering\capture[0.93]{profil.jpeg}\caption{Profil de l'utilisateur}\end{subfigure}\hfill
\begin{subfigure}[t]{0.31\textwidth}\centering\capture[0.93]{profil_preferences.jpeg}\caption{Préférences et compte}\end{subfigure}
\caption{Accueil, profil et préférences de l'application}
\label{fig:ecran_accueil}
\end{figure}

% ---------------------------------------------------------------------
\subsection{Page d'une activité}

La page d'une activité (figure~\ref{fig:ecran_activite}) affiche son logo, son type et son
code d'accès, masqué par défaut et copiable en un geste. La carte sarcelle donne le total des
dépenses ; juste en dessous, la carte « Programme » mène au programme du voyage. Les onglets
Dépenses, Membres et Photos donnent accès au détail. Le menu à trois points rassemble les
actions moins fréquentes : modifier, archiver, résumé, graphique des dépenses, statistiques
et suppression. Après l'ajout des deux dépenses du test, le total passe à 65,95 DH et chaque
dépense indique qui l'a payée.

\begin{figure}[H]
\centering
\begin{subfigure}[t]{0.31\textwidth}\centering\capture[0.93]{activite_vide.jpeg}\caption{Activité qui vient d'être créée}\end{subfigure}\hfill
\begin{subfigure}[t]{0.31\textwidth}\centering\capture[0.93]{activite_actions.jpeg}\caption{Menu des actions}\end{subfigure}\hfill
\begin{subfigure}[t]{0.31\textwidth}\centering\capture[0.93]{activite_apres_ajout.jpeg}\caption{Après l'ajout de deux dépenses}\end{subfigure}
\caption{Page d'une activité de type voyage}
\label{fig:ecran_activite}
\end{figure}

% ---------------------------------------------------------------------
\subsection{Rejoindre une activité}

Pour rejoindre un voyage, un ami saisit son code depuis l'accueil. La demande apparaît alors
chez le créateur, dans l'onglet des demandes reçues, avec le nom et le pseudonyme du
demandeur ; le créateur l'accepte ou la refuse (figure~\ref{fig:ecran_demandes}). Le
demandeur suit l'état de ses propres demandes dans l'onglet des demandes envoyées.

\begin{figure}[H]
\centering
\begin{subfigure}[t]{0.31\textwidth}\centering\capture[0.93]{demandes_recues.jpeg}\caption{Demande reçue et acceptée}\end{subfigure}\hspace{0.05\textwidth}
\begin{subfigure}[t]{0.31\textwidth}\centering\capture[0.93]{demandes_envoyees.jpeg}\caption{Demandes envoyées}\end{subfigure}
\caption{Gestion des demandes pour rejoindre une activité}
\label{fig:ecran_demandes}
\end{figure}

% ---------------------------------------------------------------------
\subsection{Ajouter une dépense en scannant un ticket}

Le bouton « Dépense » propose trois façons d'ajouter une dépense : la saisie manuelle, le
scanner de tickets et la description en une phrase (figure~\ref{fig:ecran_scan}). Après la
photo, une animation de balayage indique que l'assistant lit le ticket. Le formulaire
s'ouvre ensuite pré-rempli : sur le ticket du test, l'enseigne (« Shop Name »), le total de
47,95 et la devise ont été lus correctement. Le bandeau « Rempli à partir du ticket »
rappelle de vérifier chaque champ, et l'application signale que la date n'a pas été trouvée
sur le ticket, plutôt que d'en inventer une. Les articles lus sont reportés dans la
description, et la photo du ticket est jointe à la dépense.

\begin{figure}[H]
\centering
\begin{subfigure}[t]{0.235\textwidth}\centering\capture[0.95]{ajout_options.jpeg}\caption{Trois façons d'ajouter}\end{subfigure}\hfill
\begin{subfigure}[t]{0.235\textwidth}\centering\capture[0.95]{scan_lecture.jpeg}\caption{Lecture du ticket}\end{subfigure}\hfill
\begin{subfigure}[t]{0.235\textwidth}\centering\capture[0.95]{scan_formulaire.jpeg}\caption{Formulaire pré-rempli}\end{subfigure}\hfill
\begin{subfigure}[t]{0.235\textwidth}\centering\capture[0.95]{scan_formulaire_suite.jpeg}\caption{Répartition et articles}\end{subfigure}
\caption{Ajout d'une dépense à partir de la photo d'un ticket}
\label{fig:ecran_scan}
\end{figure}

% ---------------------------------------------------------------------
\subsection{Décrire une dépense en une phrase}

La deuxième dépense du test a été saisie en écrivant simplement « Haythame a payé le taxi
18 DH pour nous deux » (figure~\ref{fig:ecran_phrase}). L'écran rappelle les membres de
l'activité et sa devise, puis affiche « L'assistant lit ta phrase… » pendant le traitement.
Le formulaire obtenu contient le titre (Taxi), le montant (18 DH), la devise, la date du jour
(aucune date n'étant donnée dans la phrase), le payeur (Haythame Benladoul) et une
répartition à parts égales entre les deux membres. Sous chaque champ, la mention « Compris
dans ta phrase » indique son origine.

\begin{figure}[H]
\centering
\begin{subfigure}[t]{0.31\textwidth}\centering\capture[0.93]{decrire_saisie.jpeg}\caption{Phrase en cours d'analyse}\end{subfigure}\hfill
\begin{subfigure}[t]{0.31\textwidth}\centering\capture[0.93]{decrire_formulaire.jpeg}\caption{Champs compris}\end{subfigure}\hfill
\begin{subfigure}[t]{0.31\textwidth}\centering\capture[0.93]{decrire_formulaire_suite.jpeg}\caption{Payeur et répartition}\end{subfigure}
\caption{Ajout d'une dépense à partir d'une phrase libre}
\label{fig:ecran_phrase}
\end{figure}

% ---------------------------------------------------------------------
\subsection{Soldes, résumé et graphiques}

Une fois les deux dépenses enregistrées, le résumé (figure~\ref{fig:ecran_soldes}, à gauche)
indique qu'un seul remboursement suffit pour équilibrer les comptes : l'utilisateur a payé
47,95 DH pour une part de 32,98 DH, il doit donc recevoir 14,98 DH, que Haythame lui
rembourse. Le résumé peut être copié pour être envoyé au groupe. Le graphique des dépenses
présente les mêmes données sous trois angles : ce que chacun a payé (73 \% et 27 \%), ce que
chacun a consommé (50 \% chacun) et le solde qui en résulte (+14,98 DH et $-$14,98 DH).

\begin{figure}[H]
\centering
\begin{subfigure}[t]{0.235\textwidth}\centering\capture[0.95]{resume.jpeg}\caption{Résumé}\end{subfigure}\hfill
\begin{subfigure}[t]{0.235\textwidth}\centering\capture[0.95]{graphique_paye.jpeg}\caption{Payé}\end{subfigure}\hfill
\begin{subfigure}[t]{0.235\textwidth}\centering\capture[0.95]{graphique_part.jpeg}\caption{Part}\end{subfigure}\hfill
\begin{subfigure}[t]{0.235\textwidth}\centering\capture[0.95]{graphique_solde.jpeg}\caption{Solde}\end{subfigure}
\caption{Résumé des remboursements et graphiques par participant}
\label{fig:ecran_soldes}
\end{figure}

La page de statistiques (figure~\ref{fig:ecran_stats}) complète cette vue : total, nombre de
dépenses, moyenne par dépense et par participant, jour le plus dépensier, situation
personnelle de l'utilisateur, faits marquants (plus grosse dépense, plus gros payeur) et
répartition des dépenses par jour de la semaine.

\begin{figure}[H]
\centering
\begin{subfigure}[t]{0.31\textwidth}\centering\capture[0.93]{insights_1.jpeg}\caption{Indicateurs et situation}\end{subfigure}\hspace{0.05\textwidth}
\begin{subfigure}[t]{0.31\textwidth}\centering\capture[0.93]{insights_2.jpeg}\caption{Faits marquants}\end{subfigure}
\caption{Statistiques des dépenses d'une activité}
\label{fig:ecran_stats}
\end{figure}

% ---------------------------------------------------------------------
\subsection{Programme du voyage}

Pour le test, trois visites ont été ajoutées au programme : les Champs-Élysées, Notre-Dame
et la place Charles-de-Gaulle, le mercredi 14 octobre 2026 à 10 h, 12 h et 15 h
(figure~\ref{fig:ecran_programme}). Le formulaire demande un nom, un lieu, une adresse, une
date et une heure ; le lieu et l'adresse servent à placer automatiquement la visite sur la
carte. Chaque proposition peut recevoir des votes pour ou contre, affichés en direct. Le
programme se consulte en liste, en calendrier (une frise par jour, dans l'ordre des heures)
ou sur la carte Apple Plans, où chaque visite est un repère numéroté dans l'ordre
chronologique et coloré selon les votes ; un filtre permet de n'afficher qu'un jour.

\begin{figure}[H]
\centering
\begin{subfigure}[t]{0.235\textwidth}\centering\capture[0.95]{programme_ajout.jpeg}\caption{Ajout d'une visite}\end{subfigure}\hfill
\begin{subfigure}[t]{0.235\textwidth}\centering\capture[0.95]{programme_liste.jpeg}\caption{Liste et votes}\end{subfigure}\hfill
\begin{subfigure}[t]{0.235\textwidth}\centering\capture[0.95]{programme_calendrier.jpeg}\caption{Calendrier}\end{subfigure}\hfill
\begin{subfigure}[t]{0.235\textwidth}\centering\capture[0.95]{programme_carte.jpeg}\caption{Carte Apple Plans}\end{subfigure}
\caption{Programme d'un voyage : ajout, liste, calendrier et carte}
\label{fig:ecran_programme}
\end{figure}

% ---------------------------------------------------------------------
\subsection{Organisation du trajet}

Le bouton « Organiser le trajet » ouvre les options (figure~\ref{fig:ecran_trajet}) : le
mode de déplacement (à pied, à vélo, en voiture), le nombre maximal de visites par jour et
une consigne facultative. Sur le test, le trajet à pied a été demandé. Les heures des trois
visites étant déjà fixées et l'ordre déjà optimal, l'application l'indique (« Votre ordre est
déjà le meilleur ») et affiche le temps de marche total, soit 2 h 00. Pour chaque étape, elle
donne l'heure, puis le temps et la distance jusqu'à l'étape suivante (52 min et 3,9 km, puis
1 h 08 et 5,1 km), et trace l'itinéraire dans les rues sur la carte. Le bouton « Ouvrir
l'itinéraire dans Plans » lance la navigation dans l'application Plans avec toutes les
étapes, et le créateur peut enregistrer le programme proposé avec « Appliquer ce
programme ».

\begin{figure}[H]
\centering
\begin{subfigure}[t]{0.31\textwidth}\centering\capture[0.93]{trajet_options.jpeg}\caption{Options du trajet}\end{subfigure}\hspace{0.05\textwidth}
\begin{subfigure}[t]{0.31\textwidth}\centering\capture[0.93]{trajet_resultat.jpeg}\caption{Trajet proposé}\end{subfigure}
\caption{Organisation automatique du trajet d'une journée}
\label{fig:ecran_trajet}
\end{figure}

% ---------------------------------------------------------------------
\subsection{Conclusion du chapitre}

Les tests menés sur iPhone couvrent l'ensemble du parcours de l'application : création d'une
activité et arrivée d'un membre, ajout de dépenses par le scanner et par une phrase,
équilibrage des comptes, statistiques, programme partagé et organisation du trajet. Chaque
fonction d'IA a été vérifiée en conditions réelles, et chaque fois le résultat a été proposé à
l'utilisateur avant d'être enregistré. Ces tests restent menés à petite échelle (un voyage,
deux membres, quelques dépenses) : le comportement de l'application avec de grands groupes
et un usage prolongé reste à observer.
